\chapter{Lesson 7: Demonstratives and Shopping (\ja{これ}/\ja{それ}/\ja{あれ}/\ja{どれ})}

\section{Vocabulary}

\subsection{Things}
\begin{center}
\begin{tabular}{lll@{\qquad}lll}
\toprule
\ja{ほん} & hon & book & \ja{かばん} & kaban & bag \\
\ja{えんぴつ} & enpitsu & pencil & \ja{ぼうし} & bōshi & hat \\
\ja{ノート} & nōto & notebook & \ja{とけい} & tokei & watch, clock \\
\ja{かさ} & kasa & umbrella & \ja{くつ} & kutsu & shoes \\
\ja{さいふ} & saifu & wallet & \ja{Tシャツ} & tīshatsu & T-shirt \\
\ja{じしょ} & jisho & dictionary & \ja{じてんしゃ} & jitensha & bicycle \\
\ja{しんぶん} & shinbun & newspaper & \ja{ペン} & pen & pen \\
\bottomrule
\end{tabular}
\end{center}

\subsection{Places and food}
\begin{center}
\begin{tabular}{lll@{\qquad}lll}
\toprule
\ja{きっさてん} & kissaten & café & \ja{コーヒー} & kōhī & coffee \\
\ja{ぎんこう} & ginkō & bank & \ja{おちゃ} & ocha & (green) tea \\
\ja{トイレ} & toire & toilet & \ja{すし} & sushi & sushi \\
\ja{としょかん} & toshokan & library & \ja{メニュー} & menyū & menu \\
\ja{ゆうびんきょく} & yūbinkyoku & post office & \ja{とんかつ} & tonkatsu & pork cutlet \\
\ja{いちば} & ichiba & market & \ja{やさい} & yasai & vegetables \\
\bottomrule
\end{tabular}
\end{center}

\section{The ko-so-a-do system}

Japanese demonstratives come in sets that share a first syllable:
\begin{itemize}
  \item \textbf{ko-}: near the speaker;
  \item \textbf{so-}: near the listener;
  \item \textbf{a-}: far from both;
  \item \textbf{do-}: question (which?).
\end{itemize}

\begin{center}
\begin{tikzpicture}[x=1cm,y=1cm]
  % speaker and listener
  \fill[blue!10] (0,0) circle (1.3);
  \fill[orange!12] (5,0) circle (1.3);
  \node[draw, circle, fill=white, minimum size=0.8cm] (S) at (0,0) {S};
  \node[draw, circle, fill=white, minimum size=0.8cm] (L) at (5,0) {L};
  \node[font=\small] at (0,-0.75) {speaker};
  \node[font=\small] at (5,-0.75) {listener};
  \node[font=\large] at (-0.75,0.75) {\ja{これ}};
  \node[font=\large] at (5.75,0.75) {\ja{それ}};
  \node[draw, dashed, rounded corners, minimum width=2cm, minimum height=1cm, font=\large] (far) at (2.5,3) {\ja{あれ}};
  \node[font=\small] at (2.5,3.75) {far from both};
  \draw[dashed, -Stealth] (S) -- (far);
  \draw[dashed, -Stealth] (L) -- (far);
  \node[font=\large] at (2.5,-1.8) {\ja{どれ}? (which one?)};
\end{tikzpicture}
\end{center}

\begin{center}
\begin{tabular}{lllll}
\toprule
 & ko- (near me) & so- (near you) & a- (over there) & do- (which?) \\
\midrule
thing (pronoun) & \ja{これ} kore & \ja{それ} sore & \ja{あれ} are & \ja{どれ} dore \\
+ noun & \ja{この} kono & \ja{その} sono & \ja{あの} ano & \ja{どの} dono \\
place & \ja{ここ} koko & \ja{そこ} soko & \ja{あそこ} asoko & \ja{どこ} doko \\
direction (polite) & \ja{こちら} kochira & \ja{そちら} sochira & \ja{あちら} achira & \ja{どちら} dochira \\
\bottomrule
\end{tabular}
\end{center}

\begin{remark}
When speaker and listener stand side by side, the space divides into only two zones:
\ja{こ} for things near both of us, \ja{あ} for things far away. \ja{そ} is also used for
something just mentioned in conversation (``that thing you said'').
\end{remark}

\section{Grammar 1: \ja{これ}/\ja{それ}/\ja{あれ}/\ja{どれ} stand alone}

These are pronouns; they take the place of a noun and are followed by a particle.
\begin{example}
\leavevmode
\begin{itemize}
  \item \ja{これはわたしのかさです。} \emph{Kore wa watashi no kasa desu.}
        This is my umbrella.
  \item \ja{それはなんですか。} \emph{Sore wa nan desu ka.} What is that (by you)?
  \item \ja{あれはとしょかんです。} \emph{Are wa toshokan desu.} That (over there) is the
        library.
  \item \ja{やまださんのかばんはどれですか。} \emph{Yamada-san no kaban wa dore desu ka.}
        Which one is Yamada-san's bag?
\end{itemize}
\end{example}

\section{Grammar 2: \ja{この}/\ja{その}/\ja{あの}/\ja{どの} + noun}

These cannot stand alone; a noun must follow.
\begin{center}
\begin{tabular}{ll}
\toprule
Correct & Incorrect \\
\midrule
\ja{このほんはいくらですか。} How much is this book? & \ja{これほんは\ldots} \\
\ja{それはいくらですか。} How much is that? & \ja{そのはいくらですか。} \\
\bottomrule
\end{tabular}
\end{center}

\begin{example}
\leavevmode
\begin{itemize}
  \item \ja{あのひとはだれですか。} \emph{Ano hito wa dare desu ka.}
        Who is that person?
  \item \ja{このとけいはにほんのです。} \emph{Kono tokei wa Nihon no desu.}
        This watch is a Japanese one.
  \item \ja{どのじてんしゃですか。} \emph{Dono jitensha desu ka.} Which bicycle is it?
\end{itemize}
\end{example}

\section{Grammar 3: \ja{ここ}/\ja{そこ}/\ja{あそこ}/\ja{どこ} (places)}
\begin{example}
\leavevmode
\begin{itemize}
  \item[A:] \ja{すみません。ぎんこうはどこですか。} \emph{Sumimasen. Ginkō wa doko desu ka.}
        Excuse me, where is the bank?
  \item[B:] \ja{あそこです。} \emph{Asoko desu.} It's over there.
\end{itemize}
\end{example}
\ja{こちら} etc.\ are politer versions of \ja{ここ}, and also mean ``this person'' when
introducing someone: \ja{こちらはやまださんです。} ``This is Yamada-san.''

\section{Grammar 4: \ja{も} ``also''}

\ja{も} replaces \ja{は} (or \ja{が}, \ja{を}) to mean ``also/too'' when the new item
shares a property with something mentioned before.
\begin{example}
\ja{ともだちはにほんじんです。わたしもにほんじんです。}
\emph{Tomodachi wa Nihonjin desu. Watashi mo Nihonjin desu.}
My friend is Japanese. I am Japanese too.
\end{example}
\begin{important}[\ja{も} with negatives]
In negative sentences \ja{も} means ``(not) either'':
\ja{わたしもがくせいじゃないです。} ``I'm not a student either.''
\end{important}

\section{Grammar 5: noun \ja{の} (whose?) and \ja{だれの}}
\begin{example}
\ja{これはだれのさいふですか。} \emph{Kore wa dare no saifu desu ka.}
Whose wallet is this?\\
\ja{たけしさんのさいふです。} / \ja{たけしさんのです。}
It's Takeshi's (wallet).
\end{example}

\section{Shopping conversation}

Key phrases:
\begin{center}
\begin{tabular}{lll}
\toprule
Japanese & Romaji & Meaning \\
\midrule
\ja{いらっしゃいませ} & irasshaimase & welcome (clerk) \\
\ja{〜をください} & \ldots o kudasai & please give me \ldots \\
\ja{〜をおねがいします} & \ldots o onegai shimasu & \ldots, please (ordering) \\
\ja{〜をみせてください} & \ldots o misete kudasai & please show me \ldots \\
\ja{いくらですか} & ikura desu ka & how much is it? \\
\ja{じゃあ、これをください} & jā, kore o kudasai & then I'll take this \\
\ja{ひとつ}, \ja{ふたつ} & hitotsu, futatsu & one, two (items) \\
\bottomrule
\end{tabular}
\end{center}

\begin{example}[At a clothing store]
\leavevmode
\begin{itemize}
  \item[Clerk:] \ja{いらっしゃいませ。}
  \item[Me:] \ja{すみません、そのTシャツはいくらですか。}
        Excuse me, how much is that T-shirt?
  \item[Clerk:] \ja{これですか。にせんごひゃくえんです。} This one? It's ¥2,500.
  \item[Me:] \ja{じゃあ、それをください。} Then I'll take that one, please.
  \item[Clerk:] \ja{ありがとうございます。}
\end{itemize}
\end{example}

\subsection{Native counter for things}
\begin{center}
\begin{tabular}{llllll}
\toprule
1 & \ja{ひとつ} & 4 & \ja{よっつ} & 7 & \ja{ななつ} \\
2 & \ja{ふたつ} & 5 & \ja{いつつ} & 8 & \ja{やっつ} \\
3 & \ja{みっつ} & 6 & \ja{むっつ} & 9 & \ja{ここのつ} \\
 & & & & 10 & \ja{とお} \\
\bottomrule
\end{tabular}
\end{center}
``How many?'' is \ja{いくつ}. Example: \ja{おにぎりをふたつください。}
``Two rice balls, please.''

\begin{aside}
In Japan you never haggle in a normal store, and tipping is not done anywhere; trying to
tip in a restaurant often results in the staff chasing you down to return the money.
Sales tax (\ja{しょうひぜい}) was 8\% in 2017 and prices are often shown both with
and without it.
\end{aside}

\section*{Lesson 7 review}
\begin{itemize}
  \item ko (me) / so (you) / a (far) / do (which?).
  \item \ja{これ}-series stand alone; \ja{この}-series need a noun.
  \item \ja{ここ}/\ja{そこ}/\ja{あそこ}/\ja{どこ} for places; \ja{こちら} is politer.
  \item \ja{も} = also; \ja{だれの} = whose.
  \item Shopping: \ja{〜をください}, \ja{いくらですか}, native counter \ja{ひとつ}\ldots
\end{itemize}
